\section{From expert intervention to a repeatable product}

An expert user can sometimes turn a stock product's incomplete answer into a completed task. They know which context to supply, which tool to open, how to split the work, and when to check the result. The aftermarket opportunity begins when those tactics are translated into a repeatable service that an ordinary owner can use without becoming an agent operator. The harness is therefore best understood as a \emph{translation contract}: given an authorized goal, it specifies an acceptable terminal state, selects an allowed route, carries the task across changing context, and returns evidence of completion or a precise stop reason.

Figure~\ref{fig:contract} organizes this contract into four verbs. \emph{Specify} turns an open request into a task whose result can be judged. \emph{Reach} obtains the permitted context, tool, account and human checkpoint. \emph{Persist} maintains progress and exception state over a longer period than one conversational turn. \emph{Prove} checks the external state and accounts for the human work used to achieve it. This structure is intentionally about product behavior, rather than a preferred software stack. A vendor can implement it with APIs, a controlled browser, ordinary code, model calls, and human service operations in different proportions.

\begin{figure}[tbp]
\centering\begin{tikzpicture}[x=1cm,y=1cm]
\fill[ink] (0,3.6) rectangle (12.5,4.2);
\node[anchor=west,text=white,font=\small\bfseries] at (.25,3.9) {A specialist harness sells a translation contract, not a different base model};
\def\contractrow#1#2#3#4{
  \fill[teal!8] (0,#1-.35) rectangle (12.5,#1+.35);
  \node[anchor=west,text=teal,font=\small\bfseries] at (.22,#1) {#2};
  \node[anchor=west,text=ink,font=\small\bfseries] at (1.0,#1) {#3};
  \node[anchor=west,text=ink,font=\small] at (3.0,#1) {#4};
}
\contractrow{2.85}{01}{Specify}{Turn an open request into an observable goal and permitted route}
\contractrow{1.95}{02}{Reach}{Bind context, tools, accounts and human checkpoints}
\contractrow{1.05}{03}{Persist}{Carry state and exceptions across turns and waits}
\contractrow{.15}{04}{Prove}{Verify the external result and count human effort}
\end{tikzpicture}

\caption{The four obligations of an aftermarket harness. The model contributes reasoning within the contract; the specialist product makes the route repeatable and the result inspectable. The bands denote responsibilities, not a measured share of work.}
\label{fig:contract}
\end{figure}

The harness changes the effective problem presented to a model. ReAct demonstrated the usefulness of interleaving reasoning and action in its research setting \citep{react}. SWE-agent showed that a purpose-built agent-computer interface materially changes software-engineering agent performance under its benchmark setup \citep{sweagent}. Anthropic distinguishes coded workflows from agents that dynamically select steps and recommends simple composable patterns before adding complexity \citep{anthropiceffective}. These are not consumer-product A/B tests of the thesis, but they show that system design around a model is a plausible experimental variable rather than a cosmetic prompt wrapper.

A bounded harness needs a policy over interventions. On an observed stop it may ask whether the task is permitted, retrieve missing context, retry with a clearer decomposition, switch to an available authorized tool, pause for a person, or end with a stated inability. Each transition must have a budget and a success test. An unlimited loop of exhortation can turn a policy refusal into an evasion attempt, inflate cost, or produce confident text without changing the outside world. A specialist service should be able to say ``this task remains incomplete'' at a known point, with the prior attempts and open dependency preserved.

Long-running work makes persistence commercially visible. In an Anthropic engineering report, a generic long-running coding loop lost progress across context windows, while a designed harness used initialization and handoff artifacts to make sessions coherent \citep{anthropiclong}. The claim belongs to that engineering experiment. Its general lesson is a research question: when a consumer product cannot maintain state long enough to finish a workflow, can an aftermarket layer carry the state without asking the owner to reconstruct it? A durable task record is useful only if it records the goal, attempted actions, external evidence, and unresolved exceptions accurately.

This is why ``programmatic'' should be read as reproducible task handling, not merely an API call. A productized harness has an input contract, an execution policy, a verifier, an exception path, and traces that permit comparison across runs. It can be revised as the consumer product improves. Without those elements, the specialist has created a brittle demonstration, not a research result or a business. The test is whether a second owner can submit a similar permitted task and receive an accepted outcome without the original expert sitting beside the chat.

The model remains important. It can interpret varied instructions, reason over documents and interfaces, and choose among permissible actions in cases too diverse for a simple script. The harness makes that intelligence reachable and accountable. If a deterministic program reaches the same outcome more cheaply and reliably, it should replace the model for that portion of the task. Aftermarket research is not an argument for adding language-model calls to every transition; it is an argument for finding where a model plus engineered surroundings beats both a naked model and ordinary software.
